% Generated draft: requires mathematical and source review.
\documentclass[11pt]{article}
\usepackage[margin=25mm]{geometry}
\usepackage{fontspec}
\setmainfont{Latin Modern Roman}
\usepackage{amsmath,amssymb,mathtools,unicode-math}
\setmathfont{Latin Modern Math}
\usepackage{enumitem,needspace,xurl,hyperref}
\setlength{\parindent}{0pt}
\setlength{\parskip}{5pt}
\setlength{\emergencystretch}{3em}
\newcommand{\N}{\mathbb{N}}
\newcommand{\Z}{\mathbb{Z}}
\newcommand{\Q}{\mathbb{Q}}
\newcommand{\R}{\mathbb{R}}
\newcommand{\C}{\mathbb{C}}
\newcommand{\F}{\mathbb{F}}
\newcommand{\A}{\mathbb{A}}
\newcommand{\PP}{\mathbb{P}}
\newcommand{\E}{\mathbb{E}}
\newcommand{\eps}{\varepsilon}
\newcommand{\dd}{\,\mathrm{d}}
\begin{document}
\section*{UnsolvedMath ID: 10600030}
\textbf{Source problem code:} catalogue entry
\subsection{Definitions and mathematical statement}
Let \(\Phi_{\rm arr}(L)\) denote the arrow polynomial of a virtual link \(L\), and let \(H_{\rm arr}(L)\) denote the particular mod-\(2\) categorification studied in the missing reference at the start of source item 30. Write \(\operatorname{Kh}(L;\F_2)\) for the comparison Khovanov theory over \(\F_2=\Z/2\Z\). The source's arrow-state relations, homological complex, gradings, differential and comparison theory for virtual links are not supplied. In particular mod-\(2\) coefficient data must not be conflated with reducing an integer graded Euler characteristic modulo \(2\).

The source reports pairs of virtual knots \(K,K'\) for which
\[
 H_{\rm arr}(K)\not\cong H_{\rm arr}(K'),\qquad
 \operatorname{Kh}(K;\F_2)\cong\operatorname{Kh}(K';\F_2),\qquad
 \Phi_{\rm arr}(K)=\Phi_{\rm arr}(K').
\]
The isomorphisms must preserve the intended gradings, but those gradings and the example pairs are absent from the record. These are contextual observations, not an asserted theorem about every categorification called arrow homology.

Problem-list item 30 asks to understand this specific link homology through further computation and an explanation of its structure for virtual links, including the phenomenon of extra distinguishing power. It gives no proposed structure theorem, computation range or complete comparison criterion. The supplied partial-progress summary reports existing categorifications in 2009 and a homological-arrow-polynomial construction in 2022; it does not establish that all those constructions are the same theory as the one in the source or that the structural program is complete.
\subsection{Short English statement}
Investigate the source's mod-\(2\) arrow-polynomial categorification, compute more examples and explain why it can distinguish virtual knots that agree in both mod-\(2\) Khovanov homology and the arrow polynomial. Its defining complex and example pairs are missing, so the request remains tied to an unresolved source reference.
\subsection{Sources}
\begin{itemize}
\item Local UnsolvedMath registry, permanent ID \(10600030\), exact source code \texttt{catalogue entry}, virtual-knot problem-list item 30. No individual proposer or proposed year is supplied. The statement begins ``In we have studied'' with an empty reference slot; the title is truncated.
\item \emph{Unsolved Problems in Virtual Knot Theory and Combinatorial Knot Theory}, revised survey, dedicated problem-list item 30: \url{https://arxiv.org/abs/1409.2823}. Background: 60 top-level items, compound items retained intact; source-TeX extraction, accessed 2026-07-29.
\item Supplied incomplete references: ``On two categorifications of the arrow polynomial for virtual knots'' (2009), arXiv:0906.xxxx; ``The homological arrow polynomial for virtual links'' (2022), arXiv:2207.xxxx. The xxxx identifiers are placeholders, not complete references, and no authors are supplied.
\item Registry status partially\_solved and classification PARTIAL-PROGRESS report existing categorifications and ongoing structural/computational study. Source evidence NEEDS\_REVIEW presents an open problem and requires a current primary-literature audit.
\item Dataset attribution: Ulam AI / UnsolvedMath Contributors, CC BY 4.0. Source rights NEEDS\_REVIEW: publicly accessible author/publisher copy, redistribution terms requiring review.
\end{itemize}
\end{document}
